%% A paper in the shape of the CVPR author kit (article + cvpr.sty), review mode.
\documentclass[10pt,twocolumn,letterpaper]{article}

\usepackage[review]{cvpr}
\usepackage{graphicx}
\usepackage{amsmath}
\usepackage{booktabs}
\usepackage[pagebackref,breaklinks,colorlinks]{hyperref}
\usepackage[capitalize]{cleveref}

\def\confName{CVPR}
\def\confYear{2026}
\def\cvprPaperID{1234}

\begin{document}

\title{Seeing Around Corners with Commodity Cameras}

\author{First Author\\
Institution1\\
Institution1 address\\
{\tt\small firstauthor@i1.org}
\and
Second Author\\
Institution2\\
{\tt\small secondauthor@i2.org}
}

\maketitle

\begin{abstract}
We recover hidden geometry from the faint reflections a wall carries. Our method
runs on a phone and beats prior work~\cite{mildenhall2020nerf,kerbl2023gaussians,pharr2016pbrt}
by a wide margin.
\end{abstract}

\section{Introduction}
\label{sec:intro}

Non-line-of-sight imaging usually needs a laser~\cite{example2025dataset}. We show
that a wall is enough. \Cref{sec:method} describes the method and \cref{tab:main}
the results; \cref{fig:overview} gives an overview.

\section{Method}
\label{sec:method}

We model the wall as a mirror with a blur kernel $k$:
\begin{equation}
  I = k * R + n .
  \label{eq:image}
\end{equation}
Solving \cref{eq:image} for $R$ is ill-posed, so we add a learned prior.

\begin{figure}[t]
  \centering
  \includegraphics[width=0.9\linewidth]{fig/overview.png}
  \caption{Overview of the pipeline.}
  \label{fig:overview}
\end{figure}

\begin{table}[t]
  \centering
  \begin{tabular}{@{}lrr@{}}
    \toprule
    Method & PSNR $\uparrow$ & Time (s) $\downarrow$ \\
    \midrule
    Laser NLOS & 21.4 & 120 \\
    Ours & \textbf{24.1} & \textbf{0.3} \\
    \bottomrule
  \end{tabular}
  \caption{Reconstruction quality on the corner benchmark.}
  \label{tab:main}
\end{table}

{\small
\bibliographystyle{ieeenat_fullname}
\bibliography{refs}
}

\end{document}
